\section{Заключение}

В рамках работы было представлено решение для модификации трасс исполнения с целью изменения записанных образов памяти трассируемых сценариев для удовлетворения ограничений измененных настроек адресации (отличающихся от настроек трассируемой платформы). Реализованны многопроходные алгоритмы для сбора и анализа графа потока данных, которые позволяют находить локации адресов в трассе и заменять их в соответствии с новыми ограничениями. Решение позволило успешно модифицировать часть целевых трасс заказчика, которые затем прошли верификацию на функциональном симуляторе, проверку репрезентативности и были использованы для анализа производительности разрабатываемых систем с новыми настройками адресации.

Была разработана инфраструктура для описания тестовых сценариев для юнит- и модульных тестов. На основе данной инфраструктуры были написаны базовые тесты для модуля сбора ГПД. Возможно применение данной инфраструктуры для покрытия тестами других модулей.

Рассмотрена базовая методология для оценки репрезентативности (качества) выходных трасс. Представлены метрики, собранные в результате тестовых запусков решения на целевых трассах исполнения заказчика: скорость обработки (инструкций в секунду); потребление памяти (максимальный объем резидентной памяти, расход на одну инструкцию); метрики для оценки репрезентативности.

Приведены возможные направления для дальнейшего развития решения.

\clearpage
